\documentclass{beamer}
\usepackage[utf8]{inputenc}
\usepackage{graphicx}

\newtheorem{definicion}{Definición}
\newtheorem{ejemplo}{Ejemplo}

%%%%%%%%%%%%%%%%%%%%%%%%%%%%%%%%%%%%%%%%%%%%%%%%%%%%%%%%%%%%%%%%%%%%%%%%%%%%%%%
\title[Presentación con Beamer]{Práctica Beamer }
\author[SAF]{Sergio Álvarez Fernández}
\date[25-04-2014]{25 de abril de 2014}
%%%%%%%%%%%%%%%%%%%%%%%%%%%%%%%%%%%%%%%%%%%%%%%%%%%%%%%%%%%%%%%%%%%%%%%%%%%%%%%

%\usetheme{Madrid}
%\usetheme{Antibes}
%\usetheme{tree}
%\usetheme{classic}

%%%%%%%%%%%%%%%%%%%%%%%%%%%%%%%%%%%%%%%%%%%%%%%%%%%%%%%%%%%%%%%%%%%%%%%%%%%%%%%
\begin{document}
  
%++++++++++++++++++++++++++++++++++++++++++++++++++++++++++++++++++++++++++++++
\begin{frame}

  %\includegraphics[width=0.15\textwidth]{img/ullesc}
  %\hspace*{7.0cm}
  %\includegraphics[width=0.16\textwidth]{img/fmatesc}
  \titlepage

  \begin{small}
    \begin{center}
     Facultad de Matemáticas \\
     Universidad de La Laguna
    \end{center}
  \end{small}

\end{frame}
%++++++++++++++++++++++++++++++++++++++++++++++++++++++++++++++++++++++++++++++

%++++++++++++++++++++++++++++++++++++++++++++++++++++++++++++++++++++++++++++++

%++++++++++++++++++++++++++++++++++++++++++++++++++++++++++++++++++++++++++++++


\section{Primera Sección}


%++++++++++++++++++++++++++++++++++++++++++++++++++++++++++++++++++++++++++++++
\begin{frame}

\frametitle{Introducción número $\pi$}


$\pi$ (pi) es la relación entre la longitud de una circunferencia y su diámetro, en geometría euclidiana. Es un número irracional y una de las constantes matemáticas más importantes. Se emplea frecuentemente en matemáticas, física e ingeniería. El valor numérico de $\pi$, truncado a sus primeras cifras, es el siguiente:

  \begin{displaymath}  
  \pi \approx 3,14159265358979323846 \; \dots 
  \end{displaymath}



\end{frame}
%++++++++++++++++++++++++++++++++++++++++++++++++++++++++++++++++++++++++++++++

\subsection{Primera Subsección}

%++++++++++++++++++++++++++++++++++++++++++++++++++++++++++++++++++++++++++++++
\begin{frame}

\frametitle{Historia del cálculo del valor $\pi$}

La búsqueda del mayor número de decimales del número $\pi$ ha supuesto un esfuerzo constante de numerosos científicos a lo largo de la historia. Algunas aproximaciones históricas de $\pi$ son las siguientes.



\end{frame}
%++++++++++++++++++++++++++++++++++++++++++++++++++++++++++++++++++++++++++++++

\section{Segunda Sección}

%++++++++++++++++++++++++++++++++++++++++++++++++++++++++++++++++++++++++++++++
\begin{frame}
\frametitle{Mesopotamia}

Algunos matemáticos mesopotámicos empleaban, en el cálculo de segmentos, valores de $\pi$ igual a 3, alcanzando en algunos casos valores más aproximados, como el de:
  \begin{displaymath}  
  \pi \approx 3 + \frac{1}{8} = 3,125
  \end{displaymath}
  Vease \alert{Wikipedia}~\cite{latex}  
   
\end{frame}
%++++++++++++++++++++++++++++++++++++++++++++++++++++++++++++++++++++++++++++++



%++++++++++++++++++++++++++++++++++++++++++++++++++++++++++++++++++++++++++++++

%++++++++++++++++++++++++++++++++++++++++++++++++++++++++++++++++++++++++++++++

\subsection{Segunda Subsección}
%++++++++++++++++++++++++++++++++++++++++++++++++++++++++++++++++++++++++++++++
\begin{frame}
\frametitle{Época moderna}
En 1610 el matemático Ludolph van Ceulen calculó los 35 primeros decimales de $\pi$. Se dice que estaba tan orgulloso de esta hazaña que lo mandó grabar en su lápida. Los libros de matemática alemanes durante muchos años denominaron a $\pi$ como número ludolfiano. En 1665 Isaac Newton desarrolla la serie:
  \begin{displaymath}  
   \arcsin {x} = x + \frac{1}{2} \cdot \frac {x^3}{3} + \frac{1 \cdot 3}{2\cdot 4} \cdot \frac {x^5}{5} + \frac{1\cdot 3\cdot 5}{2\cdot 4\cdot 6} \cdot \frac{x^7}{7} + \ldots
  \end{displaymath}
    

Con \begin{displaymath} x = \frac {1} {2}\end{displaymath}
  obtuvo una serie para:
  \begin{displaymath}  
  \arcsin \left(\frac{1}{2}\right) = \frac{\pi}{6}
  \end{displaymath}
   

El matemático inglés John Wallis desarrolló en 1655 la conocida serie Producto de Wallis:
\begin{displaymath}  
    \frac{2}{1} \cdot \frac{2}{3} \cdot \frac{4}{3} \cdot \frac{4}{5} \cdot \frac{6}{5} \cdot \frac{6}{7} \cdot \frac{8}{7} \cdot \frac{8}{9} \cdot \dots = \frac{\pi}{2} 
\end{displaymath}

\end{frame}
\begin{frame}
\frametitle{Época moderna}

En 1699, a sugerencia de Edmond Halley, el matemático inglés Abraham Sharp (1651-1742) calculó pi con una precisión de 71 dígitos decimales usando la serie de Gregory:
\begin{displaymath} 
    \arctan (x) = x - \frac {x^3} {3} + \frac {x^5} {5} - \ldots
\end{displaymath}
Con \begin{displaymath} x = \frac {1} {\sqrt{3}}\end{displaymath} se obtiene una serie para:
\begin{displaymath} 
    \arctan \left ( \frac{1}{\sqrt{3}} \right ) = \frac{\pi}{6}
\end{displaymath}

\end{frame}
\begin{frame}
Para alcanzar la precisión obtenida, debió usar alrededor de trescientos términos en la serie. En 1720 el francés Thomas de Lagny utilizó el mismo método para obtener una aproximación de 127 dígitos (solo los primeros 112 eran correctos).

Leibniz calculó de una forma más complicada en 1682 la siguiente serie matemática que lleva su nombre:
\begin{displaymath} 
    \sum_{n=0}^{\infty} \frac{(-1)^n}{2n+1} = 1 - \frac{1}{3} + \frac{1}{5} - \dots = \frac{\pi}{4} . 
\end{displaymath}
\end{frame}

%++++++++++++++++++++++++++++++++++++++++++++++++++++++++++++++++++++++++++++++

\section{Bibliografía}
%++++++++++++++++++++++++++++++++++++++++++++++++++++++++++++++++++++++++++++++
\begin{frame}
  \frametitle{Bibliografía}

  \begin{thebibliography}{10}
   

    \beamertemplatebookbibitems
    \bibitem{latex}
     Wikipedia {\small http://es.wikipedia.org/wiki/Número\_$\pi$ }

    \beamertemplatebookbibitems
    \bibitem[URL: internet]{latex}
    Logica {\small http://www.juegosdelogica.com/numero\_pi.htm}

  \end{thebibliography}
\end{frame}

%++++++++++++++++++++++++++++++++++++++++++++++++++++++++++++++++++++++++++++++
\end{document}